% 04_implementation
\section{\bgen{РЕАЛИЗАЦИЯ}{IMPLEMENTATION}}

Тази глава описва как решенията от глава~III са изпълнени, и отделя непропорционално място
на един клас грешки. Причината е емпирична: същата грешка беше допусната три пъти на три
различни места в кода, което я прави свойство на интерфейса, а не на невниманието.

\subsection{Един интерфейс, четири механизма за вход и изход}
\label{sec:iobackends}

Четирите механизма за вход и изход стоят зад един абстрактен интерфейс от около 160 реда.
Той е умишлено тесен: колкото по-малко може да се изрази през него, толкова по-малко може да
се различава между платформите по начин, който да се появи в измерванията.

Механизмите се делят на два вида, и делението е съществено за глава~VI.

\term{Интерфейсите по готовност} съобщават, че даден дескриптор е готов, а самото системно
извикване се извършва от викащия. Такива са \code{epoll} и \code{kqueue}.
\term{Интерфейсите по завършване} приемат заявка за операция и съобщават нейното
завършване. Такива са \code{io\_uring} и IOCP.

Под Linux са налични и двата, като изборът се прави при конфигуриране, а не по платформа.
Това е и причината да съществува реализацията върху \code{epoll}: без нея механизмът за вход
и изход е напълно съвпадащ с операционната система и не може да бъде изследван като
независима променлива. Почти всички публикувани сравнения между готовност и завършване
сравняват две различни програми и следователно измерват програмите.

\subsection{Опасността при подновяване от вътрешността на \code{await\_suspend}}
\label{sec:awaiter}

Тук се описва грешка, която беше открита три пъти: веднъж наблюдавана, два пъти по
разсъждение. Тя заслужава отделен раздел, защото формата ѝ се повтаря при всеки асинхронен
интерфейс и защото поправката е промяна на интерфейса, а не на реализацията.

\subsubsection{Формата на грешката}

Асинхронната операция в модел с корутини има две части: записване на продължението, тоест
дескриптора на спряната корутина, и подаване на самата операция. Очевидната наредба е
подаване, после изчакване:

\begin{lstlisting}[language=C++,caption={Наредба, която изглежда правилна},label={lst:racebad}]
submit_operation(&op);            // операцията вече може да завърши
co_await SomeAwaiter{op};         // едва сега се записва продължението
\end{lstlisting}

Между двата реда съществува прозорец. Ако операцията завърши вътре в него, а за дейтаграма,
която вече чака в опашката на сокета, това се случва редовно, обработващият завършването
намира празно продължение, изразходва завършването и корутината не се подновява никога.
Операцията увисва.

Естествената поправка е подаването да се извърши вътре в \code{await\_suspend}, тоест след
записването на продължението. Тя затваря прозореца и въвежда втори, по-коварен.

\begin{lstlisting}[language=C++,caption={Поправка, която въвежда втора надпревара},label={lst:racebad2}]
bool await_suspend(std::coroutine_handle<> h)
{
    op.continuation = h;
    submitted = submit();     // подаването позволява завършване ВЕДНАГА
    return submitted;         // ... а дотук кадърът може вече да не съществува
}
\end{lstlisting}

Извикването \code{submit()} е точно това, което прави операцията годна за завършване.
Работна нишка може да я обработи и да поднови корутината, докато \code{submit()} още се
връща. Присвояването на \code{submitted} тогава се надпреварва с подновяването, което самото
то е предизвикало, и губи: \code{await\_resume} прочита флага, докато той още е \code{false},
и съобщава неуспешно подаване за операция, подадена съвсем успешно. По-лошото е, че
присвояването попада в кадър, който може вече да е унищожен, както и връщаната стойност
след него.

\subsubsection{Поправката е в наредбата, не в заключването}

Правилната наредба публикува оптимистичната стойност \emph{преди} подаването, докато
нишката още владее кадъра, и записва повторно само по пътя на неуспеха. Този път е безопасен
по устройство: щом нищо не е подадено, нищо не може да подновява.

\begin{lstlisting}[language=C++,caption={Наредба, която издържа},label={lst:racegood}]
bool await_suspend(std::coroutine_handle<> h)
{
    op.continuation = h;

    submitted = true;         // публикувано, докато кадърът още е наш
    if (!submit())
    {
        submitted = false;    // безопасно: нищо не е подадено, нищо не подновява
        return false;         // подновяване веднага, завършване няма да има
    }

    return true;              // литерал, а НЕ `return submitted`:
}                             // кадърът може вече да не съществува
\end{lstlisting}

Връщането на литерал вместо на член е част от поправката, а не стилистика. Всяко четене на
член след успешно подаване е достъп до евентуално освободена памет.

\subsubsection{Трите места}

Грешката беше открита в \code{io\_uring} чрез проявление: изпълнение на тестовете отказваше
приблизително веднъж на десет пъти със съобщение за неуспешно подаване на операция, която в
действителност беше подадена. Диагностиката премина през две други хипотези, всяка от които
беше поправена и не се оказа причината: липсващо заключване на опашката за подаване, която е
с един производител, и връщане от цикъла за изчакване без изчистване на опашката за
завършвания. И двете поправки са правилни сами по себе си и са отбелязани като такива, но не
обясняваха наблюдението.

Същата форма присъства в реализациите върху \code{epoll} и \code{kqueue}, където
регистрацията предхожда изчакването, а еднократната регистрация означава, че пропуснато
събитие не се повтаря никога. Там грешката е поправена \textbf{по разсъждение и не е била
възпроизведена}: двадесет изпълнения не я предизвикаха, включително с изкуствено разширен
прозорец. Причината е, че при тези реализации бързият път изобщо не регистрира, когато данни
вече чакат, а работната нишка стои в изчакване с таймаут от 100\,ms и рядко печели
надпревара, измервана в инструкции.

Това се заявява така, а не като „поправени три грешки“. Едната е наблюдавана; другите две са
затворени дупки в разсъжденията.

\subsubsection{Същата форма при отложените стойности}

Отложената стойност от раздел~\ref{sec:deferred} среща същия въпрос. Интерфейсът ѝ за
изчакване връща \emph{дали} обратното извикване е било записано, вместо да го извика:
\code{true}, ако стойността още липсва и извикването ще се случи по-късно, \code{false}, ако
стойността вече е налична, при което извикването не се прави изобщо. Изчакващият връща тази
стойност направо от \code{await\_suspend}, тоест „вече разрешено“ става „не спирай“, без
никакъв достъп до кадър, който може да е изчезнал.

Изводът, който си струва да се пренесе: при интерфейс, чиято употреба съдържа такава
надпревара, поправката е интерфейсът да прави грешката невъзможна, а не документацията да я
описва. Три повторения са достатъчно доказателство, че помненето не работи.

\subsection{Реализации на протоколите}

Трите протокола са реализирани по различни стратегии, и разликите са обосновани, а не
исторически.

\paragraph{HTTP/1.1.} Синтактичният анализ е собствен. Протоколът е достатъчно прост, за да
не оправдава зависимост, а анализаторът е и мястото, където се прилагат правилата срещу
контрабанда на заявки, тоест собствената реализация дава пряк контрол върху приемливото.

\paragraph{HTTP/2.} Рамкирането е собствено, а от \code{nghttp2} се използва само HPACK.
Разделянето е умишлено: рамкирането е механично и се проверява добре, докато HPACK носи
динамична таблица със състояние, при която грешка означава уязвимост, а не дефект.

\paragraph{HTTP/3.} Тук се използват външни библиотеки, \code{ngtcp2} и \code{nghttp3}, и
изборът им е архитектурен, а не удобство. И двете са само машини на състоянията: не
притежават нито сокет, нито цикъл на събития. Точно това прави възможен дизайнът от
глава~III. Библиотека, която управлява собствен цикъл, би наложила втори, и единственият
слушащ дескриптор щеше да отпадне, преди изобщо да бъде реализиран.

Съответно връзката по HTTP/3 наслагва три машини на състоянията: \code{ngtcp2} за
транспорта, \code{SSL} за ръкостискането, управлявано от \code{ngtcp2}, а не от BIO, и
\code{nghttp3} за рамкирането и QPACK. Заявката излиза оттам като обикновен обект заявка и
се връща като обикновен обект отговор, тоест минава през същия маршрутизатор, същата верига
от middleware и същите handler-и както останалите два протокола. Добавянето на протокол не
разклонява приложението.

\subsection{Типово безопасни параметри на маршрута}

Параметрите на маршрута пристигат като текст и почти винаги се използват като нещо друго.
Обичайният подход оставя преобразуването на handler-а: той получава асоциативен масив от
низове, изважда \code{"id"}, преобразува го към цяло число и се надява. Три неща могат да се
объркат, и трите остават невидими до изпълнение: името може да е сгрешено, преобразуването
може да не успее, а типът може да не съответства на това, което handler-ът очаква.

Регистрацията тук приема типовете като шаблонни аргументи, а изискването върху handler-а е
изразено като ограничение:

\begin{lstlisting}[language=C++,caption={Регистрация с обявени типове на параметрите},label={lst:typed}]
template <typename... Args, typename F>
    requires std::invocable<F, Args..., Request&>
void get(std::string pattern, F&& handler);

// Употреба: типовете се обявяват, handler-ът ги получава готови.
app.get<int, std::string>("/users/{id}/posts/{slug}",
    [](int id, std::string slug, Request& req) -> Task<Response>
    {
        co_return Response::ok(std::to_string(id) + ":" + slug);
    });
\end{lstlisting}

Ограничението \code{std::invocable} премества една от трите грешки от изпълнение към
компилиране: handler, чиито параметри не съответстват на обявените типове, не се компилира.
Диагностиката е на мястото на регистрацията, а не под формата на страница шаблонни
съобщения от вътрешността на библиотеката.

Преобразуването е през една точка за разширение, специализирана по тип. За целочислените
типове тя използва \code{std::from\_chars}, което не разпределя памет и не зависи от
локала. Общият случай минава през поток, но с изрична проверка, че целият вход е бил
изконсумиран: без нея \code{"12abc"} се преобразува до \code{12} и остатъкът се губи мълчаливо,
което е точно видът приемане, който по-късно се оказва разлика между два посредника.

Неуспешното преобразуване не е изключение и не е стойност по подразбиране. То е грешка,
носеща състояние 400, тоест невалиден параметър се превръща в отговор към клиента, а не в
нула, която handler-ът приема за истина.

\paragraph{Какво остава за изпълнение.} Съответствието между \emph{броя} на обявените типове
и броя на заместителите в шаблона на маршрута не се проверява при компилиране. Маршрут с два
заместителя и handler с един параметър се открива при първата заявка. Затварянето на тази
последна пролука изисква анализ на низа на шаблона по време на компилиране и е записано като
тема за отделна публикация, а не като реализирано свойство.

\subsection{\bgen{Изводи}{Summary}}

Главата описва три вида решения, и си струва да се различат, защото се защитават по различен
начин.

Първият вид е избор между налични средства с обосновка. Собствен анализатор за HTTP/1.1,
защото там живеят правилата срещу контрабанда на заявки; чужд HPACK, защото динамичната
таблица със състояние превръща дефекта в уязвимост; чужди машини на състоянията за QUIC,
защото те не притежават нито сокет, нито цикъл, а библиотека със собствен цикъл би отменила
архитектурата от глава~III, преди тя да бъде реализирана.

Вторият вид е преместване на грешка от изпълнение към компилиране. Обявените типове на
параметрите правят несъответстващия handler некомпилируем, и разделът казва изрично коя
пролука остава отворена.

Третият вид е поправка на грешка, чиято форма се повтаря. Разделът~\ref{sec:awaiter} описва
надпревара, допусната три пъти на три места, от които една беше наблюдавана, а две бяха
затворени по разсъждение и не бяха възпроизведени. Изводът е за интерфейса, а не за
конкретния код: когато обичайната употреба на интерфейса съдържа такава надпревара,
поправката е интерфейсът да прави грешката невъзможна. Три повторения са достатъчно
доказателство, че разчитането на внимание не работи.

Нито едно от твърденията в тази глава не е за производителност. Тя описва какво е
реализирано и защо; дали реализацията струва това, което се твърди, се решава в глава~VI.
